\documentclass[11pt]{article}
\usepackage{preamble}
\setlength{\parskip}{4pt}

\begin{document}

\daybanner{AP Statistics \textperiodcentered\ Unit 2 \textperiodcentered\ Topic 2.11 \textperiodcentered\ B: 2026-11-09 \textperiodcentered\ E: 2026-12-02 \textperiodcentered\ TEACHER KEY}{Check the Model}

\sectionbanner{CED TETHER}

{\small
\begin{itemize}[leftmargin=1.5em,itemsep=2pt,topsep=2pt]
  \item ENDURING UNDERSTANDING: VAR-2 The normal distribution can be used to represent some population distributions.
  \item \textbf{Skill 2.D} --- Compare distributions or relative positions of points within a distribution.
  \item \textbf{Skill 3.A} --- Determine relative frequencies, proportions, or probabilities using simulation or calculations.
  \item VAR-2.A Compare a data distribution to the normal distribution model. [Skill 2.D]
  \item VAR-2.A.1 A parameter is a numerical summary of a population. VAR-2.A.2 Some sets of data may be described as approximately normally distributed. A normal curve is mound-shaped and symmetric. The parameters of a normal distribution are the population mean, $\mu$, and the population standard deviation, $\sigma$. VAR-2.A.3 For a normal distribution, approximately 68\% of the observations are within 1 standard deviation of the mean, approximately 95\% of observations are within 2 standard deviations of the mean, and approximately 99.7\% of observations are within 3 standard deviations of the mean. This is called the empirical rule. VAR-2.A.4 Many variables can be modeled by a normal distribution.
  \item VAR-2.B Determine proportions and percentiles from a normal distribution. [Skill 3.A]
  \item VAR-2.B.1 A standardized score for a particular data value is calculated as (data value -- mean)/(standard deviation), and measures the number of standard deviations a data value falls above or below the mean. VAR-2.B.2 One example of a standardized score is a z-score, calculated as z = (x - $\mu$)/$\sigma$. A z-score measures how many standard deviations a data value is from the mean. VAR-2.B.3 Technology, such as a calculator, a standard normal table, or computer-generated output, can be used to find the proportion of data values located on a given interval of a normally distributed random variable. VAR-2.B.4 Given the area of a region under the graph of the normal distribution curve, it is possible to use technology to estimate parameters for some populations.
  \item VAR-2.C Compare measures of relative position in data sets. [Skill 2.D]
  \item VAR-2.C.1 Percentiles and z-scores may be used to compare relative positions of points within a data set or between data sets.
\end{itemize}
}
\textbf{Essential question:} When does a correct normal calculation justify a conclusion about actual data?\par
\textbf{DOK-3 skill:} 4.B \textperiodcentered\ \textbf{FRQ pattern:} {\small\texttt{is-\allowbreak{}the-\allowbreak{}normal-\allowbreak{}model-\allowbreak{}appropriate-\allowbreak{}adjudicate-\allowbreak{}model-\allowbreak{}vs-\allowbreak{}data}}\par
\textbf{DOK rationale:} Part (c) separates model arithmetic from a warranted conclusion, tests the model against shape and an observed tail, and explains the consequence of a mismatch.\par

\frameworkphaseheader{Do Now (first take) $\rightarrow$ Explore (video + follow-along) $\rightarrow$ Exit (finish + turn in)}{1 $\rightarrow$ 3}{45}{%
\begin{itemize}[leftmargin=1.5em,itemsep=2pt,topsep=2pt]
  \item Board slide up at the bell. Read Ben's and Ava's lines; do not say whether the model fits.
  \item Video follow-along from minute 5 (normal curves, empirical rule, and z-scores).
  \item Board slide back up at minute 33. Circulate; require every model judgment to point to the histogram.
  \item Collect at the door. Score part (c) only, E/P/I.
\end{itemize}
}{%
\begin{itemize}[leftmargin=1.5em,itemsep=2pt,topsep=2pt]
  \item Commit to using or rejecting the model before the video.
  \item Complete the video follow-along and check both model calculations.
  \item Finish (a)--(c); complete the exit line; turn in.
\end{itemize}
}{%
\begin{itemize}[leftmargin=1.5em,itemsep=2pt,topsep=2pt]
  \item Which claim is about a normal curve, and which claim is about the 60 actual students?
  \item How many bars continue to the right after the peak? What would symmetry require on the left?
  \item How many students does $0.62\%$ of 60 represent? How many do you actually see above 70?
  \item What must you inspect before typing values into normalcdf?
\end{itemize}
}{Solo teacher. Circulate ~5 pairs. Accept the arithmetic before pressing on the assumption; do not let `correct z-score' stand in for a model check.}

\begin{callout}{\IconWarn\ WATCH FOR}{calloutred}
\begin{itemize}[leftmargin=1.5em,itemsep=2pt,topsep=2pt]
  \item Treating the school's assertion of normality as evidence stronger than the histogram.
  \item Counting the 60--70 bar as above 70.
  \item Rejecting the model with no named feature or tail comparison.
\end{itemize}
\end{callout}

\sectionbanner{SCORING PART (c) --- E / P / I}

\textbf{Expected elements}
\begin{itemize}[leftmargin=1.5em,itemsep=2pt,topsep=2pt]
  \item \textbf{separates-calculation-from-conclusion} (required) --- Says the calculations are correct under the model but not warranted for the displayed data
  \item \textbf{cites-skew-or-shape} (required) --- Cites the peak with a long right tail or lack of mound-shaped symmetry
  \item \textbf{compares-observed-tail} (required) --- Compares 4 of 60 (6.7 percent) above 70 with the model's 0.62 percent
  \item \textbf{explains-model-check} (required) --- Names the underestimated upper tail and says to check shape before normal methods
\end{itemize}

{\small\renewcommand{\arraystretch}{1.25}
\noindent\begin{tabular}{|p{0.5in}|p{5.9in}|}
\hline
\textbf{E} & Separates correct model arithmetic from the data conclusion, cites shape and the 4-of-60 versus 0.62 percent tail discrepancy, and explains the underestimated tail and needed model check. \\ \hline
\textbf{P} & Rejects the model with only one feature, OR gives both features without explaining how correct arithmetic misleads, OR discusses model checking without comparing the tails. \\ \hline
\textbf{I} & Accepts the model because the calculations are correct, labels the histogram approximately normal, or gives no histogram evidence. \\ \hline
\end{tabular}}

\textbf{Common mistakes}
\begin{itemize}[leftmargin=1.5em,itemsep=2pt,topsep=2pt]
  \item Treating a correct $z$-score as proof that the data are normal
  \item Saying ``not normal'' without naming skew, asymmetry, or the long tail
  \item Counting the 60--70 bin as above 70 or reporting 4 of 60 as $4\%$
\end{itemize}

\clearpage
\focusbanner{TODAY'S PROBLEM --- annotated}

A school says daily app time is approximately normal with mean $40$ minutes and standard deviation $12$ minutes. The histogram shows 60 students; no recorded time was exactly 70 minutes.\par\smallskip \textbf{Ben:} ``About $16\%$ should be above 52 minutes.''\quad \textbf{Ava:} ``For 70 minutes, $z=2.5$, so less than $1\%$ should be above 70.''\par
\begin{center}
\begin{tikzpicture}
\begin{axis}[
  width=0.68\linewidth,
  height=1.84in,
  ybar interval, ymin=0, ymax=18,
  xmin=0, xmax=100, xtick={0,10,20,30,40,50,60,70,80,90,100},
  xlabel={Daily app time (minutes)}, ylabel={Number of students},
  ymajorgrids, tick label style={font=\small}, label style={font=\small},
  bar shift=0pt, x tick label as interval=false
]
\addplot[fill=calloutblue, draw=dayblue] coordinates {
  (0,3)
  (10,7)
  (20,13)
  (30,16)
  (40,9)
  (50,5)
  (60,3)
  (70,2)
  (80,1)
  (90,1)
  (100,0)
};
\end{axis}
\end{tikzpicture}
\end{center}
\begin{firsttakebox}{FIRST TAKE (5 min)}
Does the normal model fit? Point to one histogram feature.\par
\answer{Not graded. Expect trust in the stated model; preserve that tension.}
\end{firsttakebox}

\textit{Rules callout ``USING A NORMAL MODEL (keep on the page)'' is printed on the student sheet.}\par\medskip

\dokbadge{1}~\textbf{(a)} State $\mu$ and $\sigma$ for the proposed model. What two shape features should approximately normal data have?\par
\answer{$\mu=40$ min and $\sigma=12$ min. Approximately normal data should be mound-shaped and roughly symmetric.}

\dokbadge{2}~\textbf{(b)} Check Ben's and Ava's model calculations. Then find the observed fraction of these 60 students above 70 minutes.\par
\answer{52 is $1\sigma$ above 40, so Ben's normal-model result is about $16\%$. For Ava, $P(Z>2.5)\approx0.0062=0.62\%$. The histogram has $2+1+1=4$ above 70, or $4/60=6.7\%$.}

\dokbadge{3}~\textbf{(c)} Are the calculations warranted as conclusions about these students? Defend your decision with the histogram's shape and tail count. Explain how the relationship between a model and its data supports your decision, what a reader would conclude about long app times, and what must be checked before using normal methods.\par
\answer{The arithmetic is correct for the model but not warranted for these data. The histogram peaks at 30--40 and has a long right tail; 4 of 60 ($6.7\%$) exceed 70 versus the model's $0.62\%$. A calculation inherits its model assumption, so this model underestimates long app times. Check for a mound shape, symmetry, and contradictory skew, gaps, clusters, or outliers before using normal methods.}

\begin{summaryexitbox}{EXIT}
One thing the video changed about my first take: \hrulefill
\end{summaryexitbox}

\end{document}
